\documentclass{amsart}
% For dealing with references we use the comment environment
\usepackage{verbatim}
\newenvironment{reference}{\comment}{\endcomment}
%\newenvironment{reference}{}{}
\newenvironment{slogan}{\comment}{\endcomment}
\newenvironment{history}{\comment}{\endcomment}

% For commutative diagrams we use Xy-pic
\usepackage[all]{xy}

% We use 2cell for 2-commutative diagrams.
\xyoption{2cell}
\UseAllTwocells

% We use multicol for the list of chapters between chapters
\usepackage{multicol}

% This is generally recommended for better output
\usepackage{lmodern}
\usepackage[T1]{fontenc}

% For cross-file-references

% Package for hypertext links:
\usepackage{hyperref}

% For any local file, say "hello.tex" you want to link to please

% Theorem environments.
%
\theoremstyle{plain}
\newtheorem{theorem}[subsection]{Theorem}
\newtheorem{proposition}[subsection]{Proposition}
\newtheorem{lemma}[subsection]{Lemma}

\theoremstyle{definition}
\newtheorem{definition}[subsection]{Definition}
\newtheorem{example}[subsection]{Example}
\newtheorem{exercise}[subsection]{Exercise}
\newtheorem{situation}[subsection]{Situation}

\theoremstyle{remark}
\newtheorem{remark}[subsection]{Remark}
\newtheorem{remarks}[subsection]{Remarks}

\numberwithin{equation}{subsection}

% Macros
%
\def\lim{\mathop{\mathrm{lim}}\nolimits}
\def\colim{\mathop{\mathrm{colim}}\nolimits}
\def\Spec{\mathop{\mathrm{Spec}}}
\def\Hom{\mathop{\mathrm{Hom}}\nolimits}
\def\Ext{\mathop{\mathrm{Ext}}\nolimits}
\def\SheafHom{\mathop{\mathcal{H}\!\mathit{om}}\nolimits}
\def\SheafExt{\mathop{\mathcal{E}\!\mathit{xt}}\nolimits}
\def\Sch{\mathit{Sch}}
\def\Mor{\mathop{\mathrm{Mor}}\nolimits}
\def\Ob{\mathop{\mathrm{Ob}}\nolimits}
\def\Sh{\mathop{\mathit{Sh}}\nolimits}
\def\NL{\mathop{N\!L}\nolimits}
\def\CH{\mathop{\mathrm{CH}}\nolimits}
\def\proetale{{pro\text{-}\acute{e}tale}}
\def\etale{{\acute{e}tale}}
\def\QCoh{\mathit{QCoh}}
\def\Ker{\mathop{\mathrm{Ker}}}
\def\Im{\mathop{\mathrm{Im}}}
\def\Coker{\mathop{\mathrm{Coker}}}
\def\Coim{\mathop{\mathrm{Coim}}}

% Boxtimes
%
\DeclareMathSymbol{\boxtimes}{\mathbin}{AMSa}{"02}

%
% Macros for moduli stacks/spaces
%
\def\QCohstack{\mathcal{QC}\!\mathit{oh}}
\def\Cohstack{\mathcal{C}\!\mathit{oh}}
\def\Spacesstack{\mathcal{S}\!\mathit{paces}}
\def\Quotfunctor{\mathrm{Quot}}
\def\Hilbfunctor{\mathrm{Hilb}}
\def\Curvesstack{\mathcal{C}\!\mathit{urves}}
\def\Polarizedstack{\mathcal{P}\!\mathit{olarized}}
\def\Complexesstack{\mathcal{C}\!\mathit{omplexes}}
% \Pic is the operator that assigns to X its picard group, usage \Pic(X)
% \Picardstack_{X/B} denotes the Picard stack of X over B
% \Picardfunctor_{X/B} denotes the Picard functor of X over B
\def\Pic{\mathop{\mathrm{Pic}}\nolimits}
\def\Picardstack{\mathcal{P}\!\mathit{ic}}
\def\Picardfunctor{\mathrm{Pic}}
\def\Deformationcategory{\mathcal{D}\!\mathit{ef}}

\setlength{\emergencystretch}{3em}
\begin{document}
\title[Stacks Project, Tag 09J0]{712 --- Product totalization resolves complexes with exact cokernel resolutions}
\maketitle
\noindent Copyright (C) 2005--2025 Johan de Jong. Stacks Project authors. Source: \href{https://stacks.math.columbia.edu/tag/09J0}{Tag 09J0}.

\noindent Editorial difficulty note (not author text). Difficulty H2: The author constructs infinitely many compatible boundary corrections, controlling earlier corrections while extending later ones, and separately proves injectivity. This unbounded product-totalization criterion requires a substantive infinite homological argument..

\noindent Editorial provenance note (not author text). History: excerpt from revision \texttt{a0444\allowbreak 6e57e\allowbreak c1fbc\allowbreak 252a8\allowbreak 71afc\allowbreak ec775\allowbreak 2fb28\allowbreak 07b14}, homology.tex, lines 6927--7014. Reference presentation was adapted; original-author context and core support blocks were retained or added. The mathematical wording is unchanged. Licensed under GNU FDL 1.2 or later; no invariant sections or cover texts. See ../licenses/COPYING. Context blocks reproduce author definitions and supporting results; they are not additional counted problems.

\begin{lemma}
\label{lemma-good-resolution-gives-qis}
Let $M^\bullet$ be a complex of abelian groups. Let
$$
\ldots \to A_2^\bullet \to A_1^\bullet \to A_0^\bullet \to M^\bullet \to 0
$$
be an exact complex of complexes of abelian groups such that for all
$p \in \mathbf{Z}$ the complexes
$$
\ldots \to \Ker(d_{A_2^\bullet}^p) \to \Ker(d_{A_1^\bullet}^p)
\to \Ker(d_{A_0^\bullet}^p) \to \Ker(d_{M^\bullet}^p) \to 0
$$
are exact as well. Set $A^{p, q} = A_{-p}^q$ to obtain a double
complex. Then $\text{Tot}(A^{\bullet, \bullet}) \to M^\bullet$
induced by $A_0^\bullet \to M^\bullet$ is a quasi-isomorphism.
\end{lemma}

\begin{proof}
Using the short exact sequences
$0 \to \Ker(d^p_{A_n^\bullet}) \to A_n^p \to \Im(d^p_{A_n^\bullet}) \to 0$
and the assumptions we see that
$$
\ldots \to \Im(d_{A_2^\bullet}^p) \to \Im(d_{A_1^\bullet}^p)
\to \Im(d_{A_0^\bullet}^p) \to \Im(d_{M^\bullet}^p) \to 0
$$
is exact for all $p \in \mathbf{Z}$. Repeating with the exact sequences
$0 \to \Im(d^{p - 1}_{A_n^\bullet}) \to \Ker(d^p_{A_n^\bullet})
\to H^p(A_n^\bullet) \to 0$ we find that
$$
\ldots \to H^p(A_2^\bullet) \to H^p(A_1^\bullet)
\to H^p(A_0^\bullet) \to H^p(M^\bullet) \to 0
$$
is exact for all $p \in \mathbf{Z}$.

\medskip\noindent
Write $T^\bullet = \text{Tot}(A^{\bullet, \bullet})$. We will show that
$H^0(T^\bullet) \to H^0(M^\bullet)$ is an isomorphism. The same argument
works for other degrees. Let $x \in \Ker(\text{d}_{T^\bullet}^0)$ represent
an element $\xi \in H^0(T^\bullet)$.
Write $x = \sum_{i = n, \ldots, 0} x_i$ with $x_i \in A_i^i$.
Assume $n > 0$. Then $x_n$ is in the kernel of $d_{A_n^\bullet}^n$
and maps to zero in $H^n(A_{n - 1}^\bullet)$ because it maps
to an element which is the boundary of $x_{n - 1}$ up to sign.
By the first paragraph of the proof, we find that
$x_n \bmod \Im(d^{n - 1}_{A_n^\bullet})$
is in the image of $H^n(A_{n + 1}^\bullet) \to H^n(A_n^\bullet)$.
Thus we can modify $x$ by a boundary and reach the situation
where $x_n$ is a boundary. Modifying $x$ once more we see that
we may assume $x_n = 0$. By induction we see that every cohomology
class $\xi$ is represented by a cocycle $x = x_0$.
Finally, the condition on exactness of kernels tells us
two such cocycles $x_0$ and $x_0'$ are cohomologous if
and only if their image in $H^0(M^\bullet)$ are the same.
\end{proof}


\begin{lemma}
\label{lemma-good-right-resolution-gives-qis}
Let $M^\bullet$ be a complex of abelian groups. Let
$$
0 \to M^\bullet \to A_0^\bullet \to A_1^\bullet \to A_2^\bullet \to \ldots
$$
be an exact complex of complexes of abelian groups
such that for all $p \in \mathbf{Z}$ the complexes
$$
0 \to
\Coker(d_{M^\bullet}^p) \to
\Coker(d_{A_0^\bullet}^p) \to
\Coker(d_{A_1^\bullet}^p) \to
\Coker(d_{A_2^\bullet}^p) \to \ldots
$$
are exact as well. Set $A^{p, q} = A_p^q$ to obtain a double
complex. Let $\text{Tot}_\pi(A^{\bullet, \bullet})$ be the
product total complex associated to the double complex
(see proof). Then the map
$M^\bullet \to \text{Tot}_\pi(A^{\bullet, \bullet})$
induced by $M^\bullet \to A_0^\bullet$ is a quasi-isomorphism.
\end{lemma}

\begin{proof}
Abbreviating $T^\bullet = \text{Tot}_\pi(A^{\bullet, \bullet})$
we define
$$
T^n = \prod\nolimits_{p + q = n} A^{p, q} =
\prod\nolimits_{p + q = n} A_p^q
\quad\text{with}\quad
\text{d}_{T^\bullet}^n =
\prod\nolimits_{n = p + q} (f_p^q + (-1)^pd_{A_p^\bullet}^q)
$$
where $f_p^\bullet : A_p^\bullet \to A_{p + 1}^\bullet$
are the maps of complexes in the lemma.

\medskip\noindent
We will show that $H^0(M^\bullet) \to H^0(T^\bullet)$ is an isomorphism.
The same argument works for other degrees.
Let $x \in \Ker(\text{d}_{T^\bullet}^0)$ represent $\xi \in H^0(T^\bullet)$.
Write $x = (x_i)$ with $x_i \in A_i^{-i}$.
Note that $x_0$ maps to zero in $\Coker(A_1^{-1} \to A_1^0)$.
Hence we see that $x_0 = m_0 + d_{A_0^\bullet}^{-1}(y)$ for some
$m_0 \in M^0$ and $y \in A_0^{-1}$.
Then $d_{M^\bullet}(m_0) = 0$ because $\text{d}_{A_0^\bullet}(x_0) = 0$
as $\text{d}_{T^\bullet}(x) = 0$.
Thus, replacing $\xi$ by something in the image of
$H^0(M^\bullet) \to H^0(T^\bullet)$ we may assume that $x_0$
is in $\Im(d^{-1}_{A_0^\bullet})$.

\medskip\noindent
Assume $x_0 \in \Im(d^{-1}_{A_0^\bullet})$. We claim that in this
case $\xi = 0$. To prove this we find, by induction on $n$ elements
$y_0, y_1, \ldots, y_n$ with $y_i \in A_i^{-i - 1}$ such that
$x_0 = \text{d}_{A_0}^{-1}(y_0)$ and
$x_j = f_{j - 1}^{-j}(y_{j - 1}) + (-1)^j d^{-j - 1}_{A_j^\bullet}(y_j)$
for $j = 1, \ldots, n$. This is clear for $n = 0$. Proof of induction step:
suppose we have found $y_0, \ldots, y_{n - 1}$. Then
$w_n = x_n - f_{n - 1}^{-n}(y_{n - 1})$ is in the kernel of
$d^{-n}_{A_n^\bullet}$ and maps to zero in $H^{-n}(A_{n + 1}^\bullet)$
(because it maps to an element which is the boundary
of $x_{n + 1}$ up to sign). Exactly as in the proof of
Lemma \ref{lemma-good-resolution-gives-qis}
the assumptions of the lemma imply that
$$
0 \to
H^p(M^\bullet) \to
H^p(A_0^\bullet) \to
H^p(A_1^\bullet) \to
H^p(A_2^\bullet) \to \ldots
$$
is exact for all $p \in \mathbf{Z}$. Thus after changing $y_{n - 1}$
by an element in $\Ker(d^{n - 1}_{A_{n - 1}^\bullet})$ we may assume
that $w_n$ maps to zero in $H^{-n}(A_n^\bullet)$. This means we
can find $y_n$ as desired. Observe that this procedure does not
change $y_0, \ldots, y_{n - 2}$. Hence continuing ad infinitum
we find an element $y = (y_i)$ in $T^{n - 1}$ with $d_{T^\bullet}(y) = \xi$.
This shows that $H^0(M^\bullet) \to H^0(T^\bullet)$ is surjective.

\medskip\noindent
Suppose that $m_0 \in \Ker(d^0_{M^\bullet})$ maps to zero in $H^0(T^\bullet)$.
Say it maps to the differential applied to $y = (y_i) \in T^{-1}$ .
Then $y_0 \in A_0^{-1}$ maps to zero in $\Coker(d^{-2}_{A_1^\bullet})$.
By assumption this means that $y_0 \bmod \Im(d^{-2}_{A_0^\bullet})$
is the image of some $z \in M^{-1}$. It follows that
$m_0 = d^{-1}_{M^\bullet}(z)$. This proves injectivity and the proof is
complete.
\end{proof}
\end{document}
